\begin{abstract}
\noindent Researcher mobility increasingly spans more places, raising the question of how flows are redistributed as participation expands. Using author--affiliation records from 1998 to 2025, we apply the BRIDGE framework to approximately 30 million international and domestic intercity mobility events. International flows became more evenly distributed, yet 44 of the 162 countries with rising shares accounted for 85.0\% of the sum of all increases in national shares. In countries where domestic flows became more concentrated, newly recorded cities typically contributed most to rising inequality, while flows often became more evenly distributed among continuously recorded cities. The international--domestic contrast occurred in all four disciplinary domains and was strongest in social sciences. In a matched comparison holding city membership constant, the median domestic inequality trend reversed in every domain. These findings distinguish the widening geographic reach of researcher mobility from redistribution among places already participating.
\end{abstract}

